\VerdierPhysicalPageAnchor{0074}%
\VerdierSemanticAnchor{ch1-2-4-1-pi-opposite}
le foncteur opposé au foncteur défini en (2.2.3) et (2.2.4). De
même, désignons par :

\VerdierSemanticAnchor{ch1-2-4-1-eq4}
\begin{equation}
  \int_{\varphi,C^\circ}^{\Pi}\colon
  \operatorname{comp}^{n}(C^\circ)
  \longrightarrow
  \operatorname{comp}^{m}(C^\circ),
  \tag{2.4.1.4}
\end{equation}

\VerdierSemanticAnchor{ch1-2-4-1-pi-definition}
le foncteur défini en (2.2.3) et (2.2.4). On vérifie immédiatement
que le diagramme :

\VerdierSemanticAnchor{ch1-2-4-1-eq5}
\begin{equation}
  \resizebox{0.82\textwidth}{!}{%
    $\begin{tikzcd}[ampersand replacement=\&,column sep=6.5em,row sep=5.2em]
      \operatorname{comp}^{n}(C)^\circ
        \arrow[r,"\iota_n"]
        \arrow[d,"{\left(\int_{\varphi,C}^{\Sigma}\right)^\circ}"']
        \& \operatorname{comp}^{n}(C^\circ)
          \arrow[d,"{\int_{\varphi,C^\circ}^{\Pi}}"] \\
      \operatorname{comp}^{m}(C)^\circ
        \arrow[r,"\iota_m"]
        \& \operatorname{comp}^{m}(C^\circ)
    \end{tikzcd}$%
  }
  \tag{2.4.1.5}
\end{equation}

\VerdierSemanticAnchor{ch1-2-4-1-eq5-commutative}
est commutatif. En particulier, lorsque $\varphi$ est l'unique application
$[n]\longrightarrow[1]$, on a un diagramme commutatif :

\VerdierSemanticAnchor{ch1-2-4-1-eq6}
\begin{equation}
  \resizebox{0.82\textwidth}{!}{%
    $\begin{tikzcd}[ampersand replacement=\&,column sep=6.5em,row sep=5.2em]
      \operatorname{comp}^{n}(C)^\circ
        \arrow[r,"\iota_n"]
        \arrow[d,"{\left(\int_{n,C}^{\Sigma}\right)^\circ}"']
        \& \operatorname{comp}^{n}(C^\circ)
          \arrow[d,"{\int_{n,C^\circ}^{\Pi}}"] \\
      \operatorname{comp}(C)^\circ
        \arrow[r,"\iota_1"]
        \& \operatorname{comp}(C^\circ)
    \end{tikzcd}$%
  }
  \tag{2.4.1.6}
\end{equation}

\VerdierSemanticAnchor{ch1-2-4-1-twisted-isomorphism}
et on vérifie immédiatement que le couple d'isomorphismes
$(\iota_1,\iota_n)$ est un isomorphisme du $\mathbb Z[n]$-foncteur tordu
$\left(\int_{n,C}^{\Sigma}\right)^\circ$ sur le $\mathbb Z[n]$-foncteur
tordu $\int_{n,C^\circ}^{\Pi}$ (la structure graduée de
$\int_{n,C^\circ}^{\Pi}$ est décrite en (2.3.4) ; la structure graduée de
$\left(\int_{n,C}^{\Sigma}\right)^\circ$ est décrite en (2.3.4) et
(1.4.8)).

\VerdierSemanticAnchor{ch1-2-4-1-identification}
Nous identifierons dorénavant par l'isomorphisme $\iota_n$ les catégories
$\operatorname{comp}^{n}(C)^\circ$ et
$\operatorname{comp}^{n}(C^\circ)$ et les foncteurs
$\left(\int_{\varphi,C}^{\Sigma}\right)^\circ$ et
$\int_{\varphi,C^\circ}^{\Pi}$.

\VerdierSemanticAnchor{ch1-2-4-2}
\noindent\textbf{2.4.2.}\quad
Soient $C_i$, $i\in[n]$, une famille de catégories additives et
$\prod_{i\in[n]}C_i$ leur produit. On désigne par
$(X_1,X_2,\ldots,X_n)$ les objets de $\prod_{i\in[n]}C_i$ et on utilise
la notation $(X_1,X_2,\ldots,f_i,\ldots,X_n)$ pour désigner le morphisme :

\VerdierSemanticAnchor{ch1-2-4-2-component-morphism}
\begin{equation*}
  (\operatorname{id}_{X_1},\operatorname{id}_{X_2},\ldots,
  f_i,\ldots,\operatorname{id}_{X_n}).
\end{equation*}
